% Nguồn SGK Toán 9 Tập hai — Bài 23, trang 38--42:
% https://cdn3.olm.vn/upload/thuvienso/4491669493-page-38-1771844547179.png
% https://cdn3.olm.vn/upload/thuvienso/4491669493-page-39-1771844547506.png
% https://cdn3.olm.vn/upload/thuvienso/4491669493-page-40-1771844547855.png
% https://cdn3.olm.vn/upload/thuvienso/4491669493-page-41-1771844548165.png
% https://cdn3.olm.vn/upload/thuvienso/4491669493-page-42-1771844548481.png
%
% Nguồn SBT Toán 9 Tập hai — Bài 23, PDF page 29--32:
% SBT TOÁN 9 TẬP 2.pdf
% SHA-256: 5ff01fd213d1adc75f9c54b4408d6a9642a4d8038e685040efef3311a196196c

\section{Bảng tần số tương đối và biểu đồ tần số tương đối}

\begin{lesson}
    \subsection{Kiến thức trọng tâm}

    \begin{center}
    \begin{tblr}{
        width=\linewidth,
        colspec={X[2,l] X[5,l]},
        cells={valign=t},
        row{1}={font=\bfseries, halign=c}
    }
        Khái niệm, kĩ năng & Kiến thức, kĩ năng \\
        \begin{minipage}[t]{\linewidth}
        \begin{itemize}
            \item Tần số tương đối
            \item Bảng tần số tương đối
        \end{itemize}
        \end{minipage}
        &
        \begin{minipage}[t]{\linewidth}
        \begin{itemize}
            \item Thiết lập bảng tần số tương đối, biểu đồ tần số tương đối.
            \item Giải thích ý nghĩa và vai trò của tần số tương đối trong thực tiễn.
        \end{itemize}
        \end{minipage}
    \end{tblr}
    \end{center}

    Ở bài trước chúng ta đã được học các cách biểu diễn dãy dữ liệu để dễ dàng thấy được tần số xuất hiện của mỗi giá trị trong dãy. Trong bài này, chúng ta sẽ tìm hiểu các cách biểu diễn dãy dữ liệu để dễ dàng thấy được tỉ lệ xuất hiện của các giá trị trong dãy.

    \subsubsection{Bảng tần số tương đối}

    \textbf{Tần số tương đối và bảng tần số tương đối}

    \textbf{HĐ1} Có một túi kín đựng 10 quả bóng có cùng kích thước và khối lượng, mỗi quả có một trong các màu xanh, đỏ hoặc vàng. Thực hiện 30 lần lấy bóng, mỗi lần lấy 1 quả, ghi lại màu quả bóng được lấy ra sau đó trả lại bóng vào túi và trộn đều.

    \begin{enumerate}[label=\alph*)]
        \item Từ dữ liệu ghi lại, cho biết tần số xuất hiện của các quả bóng màu xanh, đỏ, vàng. Lập tỉ số giữa tần số và số lần lấy bóng.

        \answerline
        \answerline

        \item Đoán xem trong túi số lượng bóng có màu gì là ít nhất, nhiều nhất.

        \answerline
    \end{enumerate}

    \begin{keyconcept}
        Cho dãy dữ liệu $x_1,x_2,\ldots,x_n$. Tần số tương đối $f_i$ của giá trị $x_i$ là tỉ số giữa tần số của $x_i$ (gọi là $m_i$) với $n$.

        Bảng sau đây được gọi là bảng tần số tương đối.

        \begin{center}
        \begin{tblr}{
            colspec={Q[l] *{3}{Q[c]}},
            column{2-4}={mode=math}
        }
            Giá trị & x_1 & \ldots & x_k \\
            Tần số tương đối & f_1 & \ldots & f_k \\
        \end{tblr}
        \end{center}

        trong đó $n=m_1+\ldots+m_k$ và $f_1=\dfrac{m_1}{n}\cdot100(\%)$ là tần số tương đối của $x_1,\ldots$, $f_k=\dfrac{m_k}{n}\cdot100(\%)$ là tần số tương đối của $x_k$.
    \end{keyconcept}

    \textbf{Lưu ý.} Tần số tương đối còn gọi là tần suất.

    \textbf{Câu hỏi} Lập bảng tần số tương đối cho dãy dữ liệu thu được trong HĐ1.

    \answerline
    \answerline

    \textbf{Chú ý.} Người ta còn cho bảng tần số tương đối ở dạng cột: cột thứ nhất ghi các giá trị, cột thứ hai ghi tần số tương đối của các giá trị đó.

    \textbf{Ví dụ 1} Thu thập dữ liệu về chất lượng không khí tại một địa điểm trong 30 ngày mùa xuân cho kết quả như sau:

    M1, M1, M2, M2, M2, M2, M1, M2, M2, M2, M2, M2, M2, M2, M2, M4, M1, M3, M3, M3, M4, M4, M1, M1, M1, M1, M3, M3, M3, M1.

    (M1: Tốt; M2: Trung bình; M3: Kém; M4: Xấu)

    \begin{enumerate}[label=\alph*)]
        \item Lập bảng tần số tương đối cho dãy dữ liệu trên.
        \item Trong một ngày xuân, khả năng cao nhất địa điểm này có chất lượng không khí ở mức nào?
    \end{enumerate}

    \textbf{Lời giải}

    a) Tổng số ngày quan sát là $n=30$.

    Số ngày có chất lượng không khí ở các mức M1, M2, M3, M4 tương ứng là $m_1=9$, $m_2=12$, $m_3=6$, $m_4=3$.

    Do đó các tần số tương đối cho các mức M1, M2, M3, M4 lần lượt là:
    \[
        f_1=\dfrac{9}{30}\cdot100\%=30\%;\qquad
        f_2=\dfrac{12}{30}\cdot100\%=40\%;
    \]
    \[
        f_3=\dfrac{6}{30}\cdot100\%=20\%;\qquad
        f_4=\dfrac{3}{30}\cdot100\%=10\%.
    \]

    Ta có bảng tần số tương đối sau:

    \begin{center}
    \begin{tblr}{colspec={Q[l] *{4}{Q[c]}}}
        Chất lượng không khí & M1 & M2 & M3 & M4 \\
        Tần số tương đối & $30\%$ & $40\%$ & $20\%$ & $10\%$ \\
    \end{tblr}
    \end{center}

    b) Trong một ngày xuân, khả năng cao nhất là địa điểm này có chất lượng không khí ở mức M2, tức là mức Trung bình.

    \textbf{Nhận xét.} Tần số tương đối của một giá trị là ước lượng cho xác suất xuất hiện giá trị đó.

    \textbf{Luyện tập 1} Quay 50 lần một tấm bìa hình tròn được chia thành ba hình quạt với các màu xanh, đỏ, vàng. Quan sát và ghi lại mũi tên chỉ vào hình quạt có màu nào khi tấm bìa dừng lại. Kết quả thu được như sau:

    \begin{center}
    \begin{tblr}{colspec={Q[l] Q[l]}}
        Xanh & \texttt{||||/ ||||/ ||||/} \\
        Đỏ & \texttt{||||/ ||||/ ||||/ ||||/ ||||/} \\
        Vàng & \texttt{||||/ ||||/} \\
    \end{tblr}
    \end{center}

    \begin{enumerate}[label=\alph*)]
        \item Lập bảng tần số tương đối cho kết quả thu được.

        \answerline
        \answerline

        \item Ước lượng xác suất mũi tên chỉ vào hình quạt màu đỏ.

        \answerline
    \end{enumerate}

    \subsubsection{Biểu đồ tần số tương đối}

    \textbf{Tìm hiểu về biểu đồ tần số tương đối}

    \textbf{HĐ2} Có ba phương án thi đấu tại giải bóng đá khối lớp 9 của một trường như sau:

    Phương án 1: Các đội đấu vòng tròn, tính điểm;

    Phương án 2: Chia các đội thành hai bảng, mỗi bảng lấy hai đội vào trận bán kết;

    Phương án 3: Các đội bốc thăm ghép cặp, đấu loại trực tiếp.

    Ban tổ chức đã lấy phiếu khảo sát ý kiến. Kết quả được Bình và Nam biểu diễn bằng các biểu đồ như sau:

    \begin{center}
    \begin{tikzpicture}[x=1.05cm,y=0.038cm]
        \draw[->] (0,0)--(6.2,0);
        \draw[->] (0,0)--(0,90);
        \foreach \y in {0,20,40,60,80} {
            \draw (-0.08,\y)--(0.08,\y);
            \node[left] at (0,\y) {\y\%};
        }
        \foreach \x/\h/\num in {1.1/28/1,3.1/61/2,5.1/11/3} {
            \path[fill=base300] (\x-0.35,0) rectangle (\x+0.35,\h);
            \draw (\x-0.35,0) rectangle (\x+0.35,\h);
            \node[above] at (\x,\h) {\h\%};
            \node[below=2pt,align=center] at (\x,0) {Phương án\\\num};
        }
        \node[above left=1pt] at (0,90) {Tỉ lệ};
        \node[below=26pt] at (3.1,0) {Phương án thi đấu};
        \node[font=\bfseries] at (3.1,101) {Tỉ lệ ủng hộ các phương án thi đấu};
    \end{tikzpicture}

    \textit{Hình 7.9. Biểu đồ cột}
    \end{center}

    \begin{center}
    \begin{tikzpicture}
        \coordinate (O) at (0,0);
        \def\R{1.8}
        \path[fill=base300] (O)--(90:\R) arc[start angle=90,end angle=-10.8,radius=\R]--cycle;
        \path[fill=base200] (O)--(-10.8:\R) arc[start angle=-10.8,end angle=-230.4,radius=\R]--cycle;
        \path[fill=base100] (O)--(-230.4:\R) arc[start angle=-230.4,end angle=-270,radius=\R]--cycle;
        \draw (O) circle (\R);
        \draw (O)--(90:\R) (O)--(-10.8:\R) (O)--(-230.4:\R);
        \node at (40:1.0) {28\%};
        \node at (-120:0.9) {61\%};
        \node at (-250:1.2) {11\%};
        \node[right=2.3cm of O,align=left] {Phương án 1: 28\%\\Phương án 2: 61\%\\Phương án 3: 11\%};
        \node[font=\bfseries] at (0,2.35) {Tỉ lệ ủng hộ các phương án thi đấu};
    \end{tikzpicture}

    \textit{Hình 7.10. Biểu đồ hình quạt tròn}
    \end{center}

    \begin{enumerate}[label=\alph*)]
        \item Đọc và giải thích mỗi biểu đồ trên.

        \answerline
        \answerline

        \item Lập bảng tần số tương đối cho kết quả khảo sát ý kiến.

        \answerline
    \end{enumerate}

    \begin{keyconcept}
        \begin{itemize}
            \item Biểu đồ biểu diễn bảng tần số tương đối được gọi là biểu đồ tần số tương đối. Dạng thường gặp của biểu đồ tần số tương đối là biểu đồ cột và biểu đồ hình quạt tròn.
            \item Để vẽ biểu đồ hình quạt tròn ta thực hiện theo các bước sau:
        \end{itemize}

        \textbf{Bước 1.} Xác định số đo cung tương ứng của các hình quạt dùng để biểu diễn tần số tương đối của các giá trị theo công thức $360^\circ\cdot f_i$ với $i=1,\ldots,k$.

        \textbf{Bước 2.} Vẽ hình tròn và chia hình tròn thành các hình quạt có số đo cung tương ứng được xác định trong Bước 1.

        \textbf{Bước 3.} Định dạng các hình quạt tròn (thường bằng cách tô màu), ghi tần số tương đối, chú giải và tiêu đề.
    \end{keyconcept}

    \textbf{Ví dụ 2} Cho bảng tần số tương đối về loại phim yêu thích của các học sinh trong lớp 9A như sau:

    \begin{center}
    \begin{tblr}{colspec={Q[l] Q[c] Q[c] Q[c]}}
        Loại phim & Hài & Khoa học viễn tưởng & Kinh dị \\
        Tỉ lệ bạn yêu thích & $50\%$ & $37{,}5\%$ & $12{,}5\%$ \\
    \end{tblr}
    \end{center}

    Vẽ biểu đồ hình quạt tròn biểu diễn bảng tần số tương đối trên.

    \textbf{Lời giải}

    \textbf{Bước 1.} Xác định số đo cung tương ứng của các hình quạt biểu diễn các tần số tương đối cho mỗi loại phim:

    Hài: $360^\circ\cdot50\%=180^\circ$;

    Khoa học viễn tưởng: $360^\circ\cdot37{,}5\%=135^\circ$;

    Kinh dị: $360^\circ\cdot12{,}5\%=45^\circ$.

    \textbf{Bước 2.} Vẽ hình tròn và chia hình tròn thành các hình quạt (H.7.11).

    \textbf{Bước 3.} Định dạng các hình quạt tròn, ghi tỉ lệ phần trăm, chú giải và tiêu đề (H.7.12).

    \begin{center}
    \begin{tikzpicture}
        \coordinate (O) at (0,0);
        \def\R{1.8}
        \path[fill=base300] (O)--(90:\R) arc[start angle=90,end angle=-90,radius=\R]--cycle;
        \path[fill=base200] (O)--(-90:\R) arc[start angle=-90,end angle=-225,radius=\R]--cycle;
        \path[fill=base100] (O)--(-225:\R) arc[start angle=-225,end angle=-270,radius=\R]--cycle;
        \draw (O) circle (\R);
        \draw (O)--(90:\R) (O)--(-90:\R) (O)--(-225:\R);
        \node at (0:0.9) {50\%};
        \node at (-155:1.0) {37{,}5\%};
        \node at (-248:1.2) {12{,}5\%};
        \node[right=2.3cm of O,align=left] {Hài: 50\%\\Khoa học viễn tưởng: 37{,}5\%\\Kinh dị: 12{,}5\%};
        \node[font=\bfseries] at (0,2.35) {Tỉ lệ học sinh yêu thích các loại phim};
    \end{tikzpicture}

    \textit{Hình 7.12}
    \end{center}

    \textbf{Luyện tập 2} Bảng tần số tương đối sau cho biết tỉ lệ học sinh đánh giá độ khó của đề thi học kì môn Toán theo các mức độ.

    \begin{center}
    \begin{tblr}{colspec={Q[l] *{4}{Q[c]}}}
        Đánh giá & Rất khó & Khó & Trung bình & Dễ \\
        Tỉ lệ & $10\%$ & $25\%$ & $45\%$ & $20\%$ \\
    \end{tblr}
    \end{center}

    Vẽ biểu đồ hình quạt tròn biểu diễn bảng tần số tương đối này.

    \answerline
    \answerline
    \answerline

    \textbf{Tranh luận}

    Bạn Bình phát phiếu (H.7.13) lấy ý kiến bình chọn của 40 bạn trong lớp về địa điểm đi dã ngoại. Kết quả bạn Bình thu được như sau:

    \begin{center}
    \begin{tblr}{colspec={Q[l] X[c] X[c] X[c]}}
        Địa điểm & Vườn quốc gia Ba Vì & Vườn quốc gia Cát Bà & Vườn quốc gia Cúc Phương \\
        Tỉ lệ bạn bình chọn & $70\%$ & $30\%$ & $50\%$ \\
    \end{tblr}
    \end{center}

    Bình: ``Tớ sẽ dùng biểu đồ hình quạt tròn để biểu diễn bảng thống kê này.''

    Nam: ``Không được. Cậu phải dùng biểu đồ cột để biểu diễn.''

    Ý kiến của bạn thế nào?

    \begin{center}
    \fbox{\begin{minipage}{0.85\linewidth}
    \centering
    \textbf{PHIẾU BÌNH CHỌN ĐỊA ĐIỂM ĐI DÃ NGOẠI}

    (Bạn có thể lựa chọn nhiều hơn 1 địa điểm)

    A. Vườn quốc gia Ba Vì.\qquad
    B. Vườn quốc gia Cát Bà.\qquad
    C. Vườn quốc gia Cúc Phương.
    \end{minipage}}
    \end{center}

    \answerline
    \answerline
\end{lesson}

\begin{exercises}
    \subsection{Bài tập}

    \begin{questions}[resume=SGK]
        \item Lớp 9A có 40 bạn, trong đó 20 bạn mặc áo cỡ M, 13 bạn mặc áo cỡ S, 7 bạn mặc áo cỡ L. Hãy lập bảng tần số tương đối cho dữ liệu này.

        \answerline
        \answerline

        \item Biểu đồ tranh sau đây biểu diễn số lượng học sinh lớp 9B bình chọn phần mềm học trực tuyến được yêu thích nhất:

        \begin{center}
        \begin{tblr}{colspec={Q[l] Q[l]}}
            Skype & \iconUserTie\ \iconUserTie\ \iconUserTie \\
            Zoom & \iconUserTie\ \iconUserTie\ \iconUserTie\ \iconUserTie\ \iconUserTie\ \iconUserTie\ \iconUserTie\ \iconUserTie\ \iconUserTie\ \iconUserTie\ \iconUserTie \\
            Google Meet & \iconUserTie\ \iconUserTie\ \iconUserTie\ \iconUserTie\ \iconUserTie\ \iconUserTie \\
        \end{tblr}
        \end{center}

        (Mỗi \iconUserTie\ biểu diễn cho 2 học sinh)

        Lập bảng tần số tương đối cho dữ liệu được biểu diễn trong biểu đồ tranh trên.

        \answerline
        \answerline

        \item Quay 150 lần một tấm bìa hình tròn được chia thành bốn hình quạt với các màu xanh, đỏ, tím, vàng. Quan sát mũi tên chỉ vào hình quạt màu gì và ghi lại, thu được kết quả sau:

        \begin{center}
        \begin{tblr}{colspec={Q[l] *{4}{Q[c]}}}
            Màu & Xanh & Đỏ & Tím & Vàng \\
            Số lần & 60 & 30 & 40 & 20 \\
        \end{tblr}
        \end{center}

        \begin{questions}
            \item Lập bảng tần số tương đối cho dữ liệu trên.

            \answerline

            \item Ước lượng các xác suất mũi tên chỉ vào hình quạt màu xanh, màu vàng.

            \answerline

            \item Vẽ biểu đồ hình quạt tròn biểu diễn bảng tần số tương đối thu được ở câu a.

            \answerline
            \answerline
        \end{questions}

        \item Theo Tổng cục Thống kê, vào năm 2021 trong số 50,5 triệu lao động Việt Nam từ 15 tuổi trở lên có 13,9 triệu lao động đang làm việc trong lĩnh vực nông nghiệp, lâm nghiệp và thuỷ sản; 16,9 triệu lao động đang làm việc trong lĩnh vực công nghiệp và xây dựng; 19,7 triệu lao động đang làm việc trong lĩnh vực dịch vụ.

        \begin{questions}
            \item Lập bảng tần số tương đối cho dữ liệu trên.

            \answerline

            \item Vẽ biểu đồ hình quạt tròn biểu diễn bảng tần số tương đối thu được ở câu a.

            \answerline
            \answerline

            \item Tính tỉ lệ lao động không làm việc trong lĩnh vực nông nghiệp, lâm nghiệp và thuỷ sản.

            \answerline
        \end{questions}

        \item Bảng thống kê sau cho biết tỉ lệ tăng trưởng GDP năm 2022 theo khu vực kinh tế.

        \begin{center}
        \begin{tblr}{colspec={Q[l] X[c] X[c] X[c]}}
            Khu vực kinh tế & Nông nghiệp, lâm nghiệp và thủy sản & Công nghiệp và xây dựng & Dịch vụ \\
            Mức tăng trưởng & $3{,}36\%$ & $7{,}78\%$ & $9{,}99\%$ \\
        \end{tblr}
        \end{center}

        (Theo Tổng cục Thống kê)

        \begin{questions}
            \item Bảng thống kê trên có là bảng tần số tương đối hay không?

            \answerline

            \item Lựa chọn loại biểu đồ thích hợp và biểu diễn bảng thống kê trên bằng loại biểu đồ đó.

            \answerline
            \answerline
        \end{questions}
    \end{questions}
\end{exercises}

\begin{practice}
    \subsection{Luyện tập}

    \begin{questions}[resume=SBT]
        \item Một xạ thủ bắn 50 lần vào bia và ghi lại điểm của 50 lần bắn như sau:

        \begin{center}
        \begin{tblr}{colspec={Q[l] *{4}{Q[c]}}}
            Điểm & 10 & 9 & 8 & 7 \\
            Tần số & 15 & 25 & 7 & 3 \\
        \end{tblr}
        \end{center}

        \begin{questions}
            \item Lập bảng tần số tương đối cho kết quả thu được.

            \answerline

            \item Ước lượng cho xác suất xạ thủ được điểm lớn hơn hoặc bằng 9 khi bắn vào bia.

            \answerline
        \end{questions}

        \item Bạn Giang chọn một đoạn văn gồm 100 từ và đếm số chữ cái trong mỗi từ của đoạn văn này cho kết quả như trong biểu đồ sau:

        \begin{center}
        \begin{tikzpicture}[x=0.85cm,y=0.09cm]
            \draw[->] (0,0)--(9.2,0) node[below] {Số chữ cái};
            \draw[->] (0,0)--(0,39) node[above] {Tần số};
            \foreach \y in {0,5,10,15,20,25,30,35} {
                \draw (-0.08,\y)--(0.08,\y);
                \node[left] at (0,\y) {\y};
            }
            \foreach \x/\h in {1/2,2/8,3/15,4/30,5/22,6/15,7/7,8/1} {
                \draw (\x-0.25,0) rectangle (\x+0.25,\h);
                \node[above] at (\x,\h) {\h};
                \node[below] at (\x,0) {\x};
            }
            \node[font=\bfseries] at (4.5,37.5) {Số chữ cái trong 100 từ};
        \end{tikzpicture}
        \end{center}

        \begin{questions}
            \item Lập bảng tần số tương đối cho dữ liệu được biểu diễn trên biểu đồ.

            \answerline
            \answerline

            \item Tìm ước lượng cho xác suất một từ có nhiều hơn 5 chữ cái.

            \answerline
        \end{questions}

        \item Kết quả bình chọn của khán giả cho danh hiệu cầu thủ xuất sắc nhất giải bóng đá học sinh trường Trung học cơ sở Chu Văn An năm học vừa qua đối với năm cầu thủ được ban tổ chức đề cử như sau:

        \begin{center}
        \begin{tblr}{colspec={Q[l] *{5}{Q[c]}}}
            Cầu thủ & An & Tân & Việt & Bình & Lê \\
            Số lượng bình chọn & 55 & 135 & 100 & 70 & 40 \\
        \end{tblr}
        \end{center}

        \begin{questions}
            \item Lập bảng tần số tương đối cho bảng dữ liệu trên.

            \answerline

            \item Vẽ biểu đồ hình quạt tròn biểu diễn bảng tần số tương đối thu được ở câu a.

            \answerline
            \answerline
        \end{questions}

        \item Bảng sau là kết quả đánh giá của khách hàng về chất lượng phục vụ của một lái xe công nghệ, mỗi gạch biểu diễn một lần đánh giá:

        \begin{center}
        \begin{tblr}{colspec={Q[c] Q[l]}}
            $\star\star$ & \texttt{||||/ |||} \\
            $\star\star\star$ & \texttt{||||/ ||||} \\
            $\star\star\star\star$ & \texttt{||||/ ||||/ ||||/ ||} \\
            $\star\star\star\star\star$ & \texttt{||||/ |} \\
        \end{tblr}
        \end{center}

        \begin{questions}
            \item Lập bảng tần số và bảng tần số tương đối cho dữ liệu trên.

            \answerline
            \answerline

            \item Ước lượng cho xác suất một khách hàng đánh giá ở mức từ ba sao trở xuống.

            \answerline
        \end{questions}

        \item Bảng tần số tương đối sau cho biết kết quả tập luyện của một vận động viên bắn súng:

        \begin{center}
        \begin{tblr}{colspec={Q[l] *{4}{Q[c]}}}
            Điểm & 7 & 8 & 9 & 10 \\
            Tần số tương đối & $5\%$ & $15\%$ & $55\%$ & $x\%$ \\
        \end{tblr}
        \end{center}

        \begin{questions}
            \item Xác định giá trị của $x$.

            \answerline

            \item Ước lượng xác suất bắn trúng vòng 9 điểm hoặc vòng 10 điểm của vận động viên này.

            \answerline

            \item Biết rằng vận động viên đã bắn 200 viên đạn trong quá trình luyện tập. Tính số lần bắn trúng vòng 9 điểm hoặc vòng 10 điểm của vận động viên này.

            \answerline
            \answerline
        \end{questions}

        \item Cho biểu đồ hình quạt tròn sau:

        \begin{center}
        \begin{tikzpicture}
            \coordinate (O) at (0,0);
            \def\R{2.05}
            % Các phần theo thứ tự: Vật lí 23%, Hoá học 20%, Kinh tế 9%, Hoà bình 11%, Y học 24%, Văn học 13%.
            \path[fill=black!8] (O)--(90:\R) arc[start angle=90,end angle=7.2,radius=\R]--cycle;
            \path[fill=black!16] (O)--(7.2:\R) arc[start angle=7.2,end angle=-64.8,radius=\R]--cycle;
            \path[fill=black!24] (O)--(-64.8:\R) arc[start angle=-64.8,end angle=-97.2,radius=\R]--cycle;
            \path[fill=black!32] (O)--(-97.2:\R) arc[start angle=-97.2,end angle=-136.8,radius=\R]--cycle;
            \path[fill=black!48] (O)--(-136.8:\R) arc[start angle=-136.8,end angle=-223.2,radius=\R]--cycle;
            \path[fill=black!64] (O)--(-223.2:\R) arc[start angle=-223.2,end angle=-270,radius=\R]--cycle;
            \draw (O) circle (\R);
            \draw (O)--(90:\R) (O)--(7.2:\R) (O)--(-64.8:\R) (O)--(-97.2:\R) (O)--(-136.8:\R) (O)--(-223.2:\R);
            \node[font=\bfseries,align=center] at (0,2.75) {Tỉ lệ người đạt giải Nobel theo các lĩnh vực\\tính đến năm 2020};
            \node[align=center] at (2.8,1.1) {Vật lí\\23\%};
            \node[align=center] at (2.9,-0.25) {Hoá học\\20\%};
            \node[align=center] at (2.3,-1.55) {Kinh tế\\9\%};
            \node[align=center] at (-2.55,-1.55) {Hoà bình\\11\%};
            \node[align=center] at (-2.75,-0.15) {Y học\\24\%};
            \node[align=center] at (-2.5,1.35) {Văn học\\13\%};
        \end{tikzpicture}
        \end{center}

        \begin{questions}
            \item Giải thích các số liệu được biểu diễn trên biểu đồ.

            \answerline
            \answerline

            \item Lập bảng tần số tương đối cho dữ liệu được biểu diễn trên biểu đồ.

            \answerline
        \end{questions}

        \item Cho bảng thống kê sau về số lượng học sinh tại một trường tham gia các câu lạc bộ (CLB):

        \begin{center}
        \begin{tblr}{colspec={Q[l] *{4}{Q[c]}}}
            & CLB tiếng Anh & CLB Toán & CLB Khoa học & Tổng \\
            Nam & 40 & 60 & 50 & 150 \\
            Nữ & 70 & 30 & 65 & 165 \\
            Tổng & 110 & 90 & 115 & 315 \\
        \end{tblr}
        \end{center}

        \begin{questions}
            \item Lập các bảng tần số tương đối biểu diễn tỉ lệ học sinh nam, nữ tham gia các câu lạc bộ.

            \answerline
            \answerline

            \item Vẽ các biểu đồ hình quạt tròn biểu diễn các bảng tần số tương đối thu được ở câu a.

            \answerline
            \answerline
        \end{questions}

        \item Với số liệu cho trong bài tập 7.15, hãy:
        \begin{questions}
            \item Lập bảng tần số tương đối cho tỉ lệ học sinh tham gia các câu lạc bộ.

            \answerline

            \item Vẽ biểu đồ cột biểu diễn bảng tần số tương đối thu được ở câu a.

            \answerline
            \answerline
        \end{questions}
    \end{questions}
\end{practice}
